\documentclass[11pt,a4paper,openany]{book}
\usepackage[margin=24mm]{geometry}
\usepackage{fontspec}
\IfFontExistsTF{Noto Serif}{\setmainfont{Noto Serif}}{\setmainfont{DejaVu Serif}}
\usepackage{amsmath,amssymb}
\usepackage{tikz}
\usepackage{enumitem}
\usepackage{microtype}
\usepackage[hidelinks]{hyperref}
\setlength{\parindent}{0pt}
\setlength{\parskip}{5pt}
\setlist[enumerate]{leftmargin=*,itemsep=9pt}

\renewcommand{\contentsname}{Mục lục}
\begin{document}
\title{V-TALENT 2045\\ PHY102-A --- Chuyển động, đồ thị và vận tốc tương đối\\\large Level 1 · Cơ học nền tảng · Lớp 8–9}
\author{Dr. Dang Minh Tuan -- TS. Đặng Minh Tuấn}
\date{2026 · v1.0}
\maketitle
\tableofcontents
\chapter*{Mục tiêu năng lực}\addcontentsline{toc}{chapter}{Mục tiêu năng lực}
Phân biệt quãng đường và độ dời; đọc độ dốc và diện tích; xây dựng mô hình chuyển động từng đoạn; giải bài gặp nhau, đuổi nhau; giải thích giả thiết và kiểm tra kết quả.

Tập chuyên đề đầu tiên, bản thảo 1.0. Bài tập tự biên soạn; chỉ mục đề ngoài dùng để đọc thêm, không sao chép. Cần thẩm định độc lập trước khi dùng đánh giá chính thức.
\chapter{Mô tả chuyển động: xây dựng mô hình thay vì học thuộc}
\section{Lý thuyết và suy luận}
Tọa độ $x$ là đại lượng có dấu trong hệ quy chiếu đã chọn. Độ dời $\Delta x=x_2-x_1$; quãng đường $s$ là tổng chiều dài đường đi và không âm. Trong một chiều, $s\geq|\Delta x|$, dấu bằng khi không đổi chiều (không kể khoảng đứng yên). Phải chọn hệ quy chiếu và chiều dương trước khi gán dấu. Vật có thể ở tọa độ âm nhưng vẫn chuyển động theo chiều dương.

Tốc độ trung bình $\bar u=s/\Delta t$; vận tốc trung bình $\bar v=\Delta x/\Delta t$. Nếu hai chặng có thời gian $t_1,t_2$ và tốc độ $u_1,u_2$ thì $\bar u=(u_1t_1+u_2t_2)/(t_1+t_2)$. Với hai quãng đường bằng nhau, $\bar u=2u_1u_2/(u_1+u_2)$; với hai thời gian bằng nhau, đó là trung bình cộng. Hai công thức xuất phát từ điều kiện lấy trọng số khác nhau.

Chuyển động thẳng đều tuân theo $x=x_0+vt$, vì độ dời bằng nhau trong các khoảng thời gian bằng nhau. Với gia tốc không đổi, đồ thị vận tốc--thời gian là đường thẳng: $v=v_0+at$. Diện tích đại số hình thang cho $\Delta x=(v_0+v)t/2=v_0t+at^2/2$. Đây là nội dung mở rộng Level 1, suy ra bằng hình học không cần giải tích. Khi gia tốc thay đổi, không được áp dụng các công thức này cho cả quá trình.

\section{Ví dụ minh họa}
\subsection*{Ví dụ minh họa 1.1}
Một học sinh đi 120 m về phía đông trong 60 s, chờ 30 s, rồi đi 80 m về phía tây trong 40 s. Tính quãng đường, độ dời, tốc độ và vận tốc trung bình.

\textbf{Lời giải.} Chọn chiều đông dương. Thời gian toàn bộ kể cả lúc chờ: $t=130\,\mathrm{s}$. Quãng đường $120+80=200\,\mathrm{m}$, độ dời $\Delta x=120-80=40\,\mathrm{m}$. Suy ra $\bar u=200/130=1{,}54\,\mathrm{m\,s^{-1}}$ và $\bar v=40/130=0{,}308\,\mathrm{m\,s^{-1}}$ hướng đông. Hai đại lượng cùng đơn vị nhưng khác ý nghĩa.
\subsection*{Ví dụ minh họa 1.2}
Người đi xe đạp đi hai quãng đường bằng nhau với tốc độ 12 km/h và 18 km/h. Suy ra tốc độ trung bình và so sánh với trung bình cộng.

\textbf{Lời giải.} Gọi mỗi quãng đường là $d$ km. Khi đó $t=d/12+d/18=5d/36$ giờ và $\bar u=2d/(5d/36)=14{,}4$ km/h. Trung bình cộng là 15 km/h. Chặng chậm kéo dài hơn nên có trọng số thời gian lớn hơn; kết quả nằm giữa 12 và 18.
\section{Luyện tập cuối chuyên đề}
\begin{enumerate}
\item \textbf{PHY102-A-1.1}. Đi 90 m về bắc trong 45 s và 30 m về nam trong 15 s. Xác định bốn đại lượng chuyển động đã học.
\item \textbf{PHY102-A-1.2}. Xe buýt đi 40 km với 40 km/h, dừng 15 phút, rồi đi 60 km với 60 km/h. Tính tốc độ trung bình kể cả lúc dừng.
\item \textbf{PHY102-A-1.3}. Với hai quãng đường bằng nhau, tốc độ trung bình 24 km/h. Một chặng đi 20 km/h. Tìm tốc độ chặng kia.
\item \textbf{PHY102-A-1.4}. Vật bắt đầu tại x = -5 m, chuyển động đều với +3 m/s trong 4 s rồi -2 m/s trong 5 s. Tìm tọa độ cuối và quãng đường.
\end{enumerate}
\chapter{Đồ thị, bài toán gặp nhau và hệ quy chiếu chuyển động}
\section{Lý thuyết và suy luận}
Đồ thị tọa độ--thời gian mã hóa chuyển động: độ dốc $v=\Delta x/\Delta t$ đối với đoạn thẳng, đoạn ngang là đứng yên, dốc âm là chuyển động ngược chiều dương. Hai đồ thị giao nhau biểu thị cùng tọa độ, cùng thời điểm, không nhất thiết va chạm nếu ở hai làn khác nhau. Trên đồ thị vận tốc--thời gian, diện tích đại số là độ dời; tính quãng đường bằng cách chia tại điểm đổi dấu rồi cộng giá trị tuyệt đối. Đồ thị tọa độ dốc lớn không đồng nghĩa gia tốc lớn.

Với chuyển động đều, dùng một mốc thời gian chung: $x_A=x_{A0}+v_A t$, $x_B=x_{B0}+v_B t$. Khoảng cách đại số $D=x_B-x_A=D_0+(v_B-v_A)t$. Đặt $D=0$ cho thời điểm gặp nếu nghiệm không âm và nằm trong khoảng mô hình. Khi B xuất phát muộn, $x_B=x_{B0}+v_B(t-t_d)$ chỉ đúng với $t\ge t_d$; trước đó B ở $x_{B0}$.

Vận tốc tương đối suy từ hiệu tọa độ: $v_{A/B}=v_A-v_B$. Trong hai chiều đây là hệ thức vector. Với người bơi trong dòng đều, $\vec v_{\rm đất}=\vec v_{\rm nước}+\vec v_{\rm dòng}$. Hướng bơi vuông góc bờ cho thời gian qua sông nhỏ nhất ở tốc độ bơi cố định, nhưng bị trôi. Có thể hướng chếch thượng lưu để triệt độ trôi nếu tốc độ bơi lớn hơn tốc độ dòng.
\begin{center}
\begin{tikzpicture}[x=.85cm,y=.52cm]
\draw[->] (0,0)--(7,0) node[right] {$t$};
\draw[->] (0,-3)--(0,5) node[above] {$v$};
\draw[thick] (0,-2)--(6,4);
\draw[dashed] (2,0)--(2,-2);
\draw[dashed] (6,0)--(6,4);
\node[below] at (2,0) {$t_0$};\node[below] at (6,0) {$t_1$};
\end{tikzpicture}
\end{center}
\section{Ví dụ minh họa}
\subsection*{Ví dụ minh họa 2.1}
A xuất phát lúc 0 giờ với 5 km/h. B xuất phát cùng nơi sau 0,5 giờ với 8 km/h. Khi nào và ở đâu B đuổi kịp A?

\textbf{Lời giải.} Lúc $t=0{,}5$ giờ, A dẫn 2,5 km. Tốc độ thu hẹp khoảng cách $8-5=3$ km/h. Thời gian còn lại $2{,}5/3=5/6$ giờ; thời gian chung $4/3$ giờ. Vị trí $5(4/3)=20/3$ km. Tương đương giải $5t=8(t-0{,}5)$ và kiểm tra $t\ge0{,}5$.
\subsection*{Ví dụ minh họa 2.2}
Sông rộng 120 m. Người bơi với tốc độ 1,5 m/s so với nước theo hướng vuông góc bờ. Dòng chảy 0,9 m/s. Tính thời gian qua sông và độ trôi. Sau đó tính thời gian nếu đến đúng điểm đối diện.

\textbf{Lời giải.} Thành phần qua sông 1,5 m/s nên $t=120/1{,}5=80$ s và độ trôi $0{,}9(80)=72$ m. Muốn triệt độ trôi, thành phần hướng thượng lưu cần bằng 0,9 m/s; theo Pythagoras thành phần qua sông $\sqrt{1{,}5^2-0{,}9^2}=1{,}2$ m/s. Do đó $t=100$ s. Đi thẳng đến đối diện lâu hơn vì một phần vận tốc bơi dùng để bù dòng.
\section{Nhiệm vụ thực nghiệm}
Quay xe đồ chơi cạnh thước với khoảng thời gian giữa khung hình đã biết. Lấy ít nhất sáu cặp (t,x), vẽ x theo t, ước lượng độ dốc từ hai điểm cách xa rồi so sánh các đoạn. Báo cáo hiệu chuẩn pixel–chiều dài và bất định thời gian. Nếu độ dốc khác nhau vượt bất định, bác bỏ mô hình đều thay vì lấy trung bình để che khác biệt.
\section{Luyện tập cuối chuyên đề}
\begin{enumerate}
\item \textbf{PHY102-A-2.1}. A và B cách nhau 42 km, đi về phía nhau với 12 và 9 km/h. Tính thời gian và vị trí gặp tính từ A.
\item \textbf{PHY102-A-2.2}. Xe A đi 60 km/h. Xe B xuất phát cùng nơi sau 20 phút với 75 km/h. Tính thời gian đuổi kịp kể từ lúc A xuất phát.
\item \textbf{PHY102-A-2.3}. Vận tốc tăng đều từ -4 đến +8 m/s trong 6 s. Tính gia tốc, độ dời và quãng đường.
\item \textbf{PHY102-A-2.4}. Thuyền có tốc độ 5 m/s so với nước. Sông rộng 240 m chảy với 3 m/s. Tìm thời gian qua sông nhỏ nhất, độ trôi và thời gian đến đúng điểm đối diện.
\end{enumerate}
\chapter*{Tài liệu và bài đọc mở rộng}\addcontentsline{toc}{chapter}{Tài liệu và bài đọc mở rộng}
OpenStax Physics; PhET; Tracker; kho SJPO. Bài VPhO là phần mở rộng Level 2, không bắt buộc ở cấp độ này.

\url{https://diavn.github.io/truongchuyen/physics-vpho.html}

\url{https://hyperphysics.gsu.edu/hbase/hph.html}

\url{https://openstax.org/subjects/science}
\appendix
\chapter{Lời giải đầy đủ phần luyện tập}
\section{Mô tả chuyển động: xây dựng mô hình thay vì học thuộc}
\begin{enumerate}
\item \textbf{PHY102-A-1.1}. Chọn bắc dương: $s=120$ m, $\Delta x=60$ m, $t=60$ s. Do đó $\bar u=2$ m/s, $\bar v=1$ m/s về bắc.
\item \textbf{PHY102-A-1.2}. Mỗi chặng mất 1 giờ; dừng 0,25 giờ. Vậy $t=2{,}25$ giờ, $s=100$ km, $\bar u=100/2{,}25=44{,}44$ km/h. Nếu bỏ lúc dừng sẽ ra 50 km/h, là một bài toán khác.
\item \textbf{PHY102-A-1.3}. Dùng $2/24=1/20+1/u$. Suy ra $1/u=1/12-1/20=1/30$, nên $u=30$ km/h. Kiểm tra: $2(20)(30)/(20+30)=24$.
\item \textbf{PHY102-A-1.4}. Độ dời chặng đầu $+12$ m cho $x=7$ m. Chặng sau $-10$ m cho $x_f=-3$ m. Quãng đường $12+10=22$ m, không phải $|-3-(-5)|=2$ m.
\end{enumerate}
\section{Đồ thị, bài toán gặp nhau và hệ quy chiếu chuyển động}
\begin{enumerate}
\item \textbf{PHY102-A-2.1}. Chọn vị trí đầu của A làm gốc. $x_A=12t$, $x_B=42-9t$. Cho bằng nhau: $21t=42$, $t=2$ giờ. $x=24$ km tính từ A.
\item \textbf{PHY102-A-2.2}. Trễ $1/3$ giờ, khoảng cách ban đầu $20$ km, tốc độ tương đối 15 km/h. B mất $20/15=4/3$ giờ; tổng thời gian $5/3$ giờ = 100 phút.
\item \textbf{PHY102-A-2.3}. $a=(8-(-4))/6=2$ m/s\textsuperscript{2}. Vận tốc bằng 0 ở $t=2$ s. Tam giác âm có diện tích $-4(2)/2=-4$ m; tam giác dương $8(4)/2=16$ m. Độ dời $12$ m; quãng đường $20$ m.
\item \textbf{PHY102-A-2.4}. Thời gian nhỏ nhất $240/5=48$ s, độ trôi $3(48)=144$ m. Muốn bù dòng, thành phần ngược dòng 3 m/s và thành phần ngang sông $\sqrt{25-9}=4$ m/s. Thời gian đến đối diện $240/4=60$ s. Kiểm tra tính khả thi: $5>3$.
\end{enumerate}
\end{document}
